
\documentclass{article}
\usepackage{colortbl}
\usepackage{makecell}
\usepackage{multirow}
\usepackage{supertabular}

\begin{document}

\newcounter{utterance}

\centering \large Interaction Transcript for game `cladder', experiment `full\_v1.5\_default', episode 8348 with llama{-}3.1{-}8b{-}instruct.
\vspace{24pt}

{ \footnotesize  \setcounter{utterance}{1}
\setlength{\tabcolsep}{0pt}
\begin{supertabular}{c@{$\;$}|p{.15\linewidth}@{}p{.15\linewidth}p{.15\linewidth}p{.15\linewidth}p{.15\linewidth}p{.15\linewidth}}
    \# & \multicolumn{2}{c}{Player} && \multicolumn{2}{c}{Game Master} \\
    \hline

    \theutterance \stepcounter{utterance}  
    & & & \multicolumn{4}{p{0.6\linewidth}}{
        \cellcolor[rgb]{0.9,0.9,0.9}{
            \makecell[{{p{\linewidth}}}]{
                \texttt{\tiny{[P1$\langle$GM]}}
                \texttt{You are an expert in causal inference. The following question is not a typical commonsense query, but rather a meticulously designed question created by a professor specializing in causal inference, intended to assess the students' mastery of the course content.} \\
\texttt{Imagine a self{-}contained, hypothetical world with only the following conditions, and without any unmentioned factors or causal relationships: Yomx has a direct effect on cwoi and gwet. Gwet has a direct effect on xevu. Cwoi has a direct effect on xevu. The overall probability of yomx is 22\%. The probability of not yomx and xevu is 48\%. The probability of yomx and xevu is 7\%. Is the chance of xevu larger when observing yomx?} \\
\texttt{Start your answer with 'Yes' or 'No', followed by additional reasoning or evidence to support your explanation.} \\
            }
        }
    }
     \\ \\

    \theutterance \stepcounter{utterance}  
    & \multicolumn{4}{p{0.6\linewidth}}{
        \cellcolor[rgb]{0.561,0.737,0.561}{
            \makecell[{{p{\linewidth}}}]{
                \texttt{\tiny{[P1$\rangle$GM]}}
                \texttt{No.} \\
\\ 
\texttt{To determine if the chance of xevu is larger when observing yomx, we need to calculate the conditional probability of xevu given yomx (P(xevu|yomx)) and compare it to the unconditional probability of xevu (P(xevu)).} \\
\\ 
\texttt{We can use the law of total probability to calculate P(xevu|yomx) and P(xevu):} \\
\\ 
\texttt{P(xevu|yomx) = P(yomx and xevu) / P(yomx)} \\
\texttt{= 7\% / 22\%} \\
\texttt{= 0.3182 (or approximately 31.82\%)} \\
\\ 
\texttt{P(xevu) = P(not yomx and xevu) + P(yomx and xevu)} \\
\texttt{= 48\% (since the probability of not yomx and xevu is 48\%, and the probability of xevu is contained within this group)} \\
\texttt{= 0.48} \\
\\ 
\texttt{To determine if the chance of xevu is larger when observing yomx, we need to compare P(xevu|yomx) to P(xevu). Since P(xevu|yomx) = 0.3182 and P(xevu) = 0.48, we can see that P(xevu) is greater than P(xevu|yomx). This indicates that the chance of xevu is not larger when observing yomx.} \\
\\ 
\texttt{Therefore, the answer is No.} \\
            }
        }
    }
    & & \\ \\

    \theutterance \stepcounter{utterance}  
    & & & \multicolumn{2}{p{0.3\linewidth}}{
        \cellcolor[rgb]{0.9,0.9,0.9}{
            \makecell[{{p{\linewidth}}}]{
                \texttt{\tiny{[GM$|$GM]}}
                \texttt{No} \\
            }
        }
    }
    & & \\ \\

    \theutterance \stepcounter{utterance}  
    & & & \multicolumn{2}{p{0.3\linewidth}}{
        \cellcolor[rgb]{0.9,0.9,0.9}{
            \makecell[{{p{\linewidth}}}]{
                \texttt{\tiny{[GM$|$GM]}}
                \texttt{no} \\
            }
        }
    }
    & & \\ \\

    \theutterance \stepcounter{utterance}  
    & & & \multicolumn{2}{p{0.3\linewidth}}{
        \cellcolor[rgb]{0.9,0.9,0.9}{
            \makecell[{{p{\linewidth}}}]{
                \texttt{\tiny{[GM$|$GM]}}
                \texttt{game\_result = WIN} \\
            }
        }
    }
    & & \\ \\

\end{supertabular}
}

\end{document}
